\section{Verdikt}
\label{sec:verdikt}

Dieser Teil zieht die Schlussfolgerung aus den Befunden der
Teile~\ref{sec:organisch} bis \ref{sec:urteil}. Er f\"uhrt keine neue Zahl
ein; jede genannte Gr\"o{\ss}e steht bereits oben.

\vd{\textbf{Der rote Text ist eine Empfehlung, kein Beleg.}} Er ist aus den
vorstehenden Befunden unter den in Teil~\ref{sec:flow} erkl\"arten Annahmen
abgeleitet. Schwarz bleibt, was belegt ist; blau sind Einsch\"atzungen im
laufenden Text; rot ist das Votum. Die Entscheidung trifft der Verfasser.

\subsection{Votum}

\vd{\textbf{Nicht kaufen zum aktuellen Kurs von \USDje{\FlowEinstiegskurs}.
Beobachten und auf einen der beiden Ausl\"oser in Abschnitt~\ref{sec:ausloeser}
warten.} Kein Verkaufsvotum: Wer die Aktie h\"alt, hat keinen aus dieser
Analyse ableitbaren Grund zu verkaufen, denn die Verschuldung ist tragbar und
das Gesch\"aft erwirtschaftet Bargeld. Wer sie nicht h\"alt, hat keinen Grund
einzusteigen, denn er bezahlt den fairen Wert ohne Sicherheitsabstand.}

\subsection{Begr\"undung aus den Befunden}

\vd{Drei Befunde tragen das Votum, und sie zeigen in verschiedene Richtungen.

\emph{Gegen einen Kauf spricht der Preis.} Der interne Zinsfu{\ss} betr\"agt
im mittleren Szenario \Pct{\FlowIRRBFuenf} \"uber f\"unf und
\Pct{\FlowIRRBFuenfzehn} \"uber f\"unfzehn Jahre
(Tabelle~\ref{tab:flowirr}), der Ma{\ss}stab der eigenen Eigenkapitalrendite
\Pct{\FlowHuerde}. Auf dem l\"angsten Horizont wird er also knapp verfehlt.
Ein Kauf zum heutigen Kurs bringt ungef\"ahr das, was das Unternehmen auf sein
eigenes Kapital verdient, und keinen Aufschlag daf\"ur, das Risiko zu tragen.

\emph{Gegen einen Kauf spricht auch der Zeitpunkt.} Der Fr\"uhindikator, der
den Kursverlauf der letzten zwei Jahre erkl\"art (Teil~\ref{sec:kurs}), ist
negativ geworden: \Pct{\HJOrganisch} im ersten Halbjahr 2026 nach
\Pct{\OrganischWachstumVJ} im Gesch\"aftsjahr 2024. Der berichtete
Umsatzzuwachs von \Pct{\HJUmsatzDelta} ist ein Konsolidierungseffekt und
neutralisiert sich ab August 2026 von selbst. Es gibt in den n\"achsten vier
bis sechs Quartalen keinen belegbaren Ausl\"oser f\"ur eine Neubewertung.

\emph{F\"ur ein Beobachten statt eines Meidens spricht die Bilanz.} Die
Verschuldung, die den Kurs gedr\"uckt hat, ist nicht das Problem: Der
Verschuldungsgrad f\"allt in \emph{jedem} der drei gerechneten Szenarien, im
ung\"unstigsten von \Faktor{\FlowGradStart} auf
\Faktor{\FlowGradSzenarioA} (Tabelle~\ref{tab:flowszenarien}). Ein Makler
bindet kaum Kapital: \MioUSD{\FlowSachinvest} Sachinvestitionen gegen
\MioUSD{\FlowEbitdacEnde} EBITDAC. Das Gesch\"aft ist gesund, der Preis ist es
nicht.}

\subsection{Einstiegsschwelle}

Tabelle~\ref{tab:schwelle} nennt den Kurs, bei dem der interne Zinsfu{\ss}
den Ma{\ss}stab von \Pct{\FlowHuerde} genau trifft. Das ist die Zahl, die eine
Anlageentscheidung braucht: nicht ob, sondern bis zu welchem Preis.

\begin{table}[htbp]
\centering
\caption{Einstiegskurs in USD, bei dem der interne Zinsfu{\ss} die H\"urde von
\Pct{\FlowHuerde} erreicht. Ausstiegsvielfaches unver\"andert bei
\num{\FlowVielfaches}. Aktueller Kurs: \USDje{\FlowEinstiegskurs}.}
\label{tab:schwelle}
\begin{tabular}{lrrrr}
\toprule
Szenario & Wachstum in \% & 5 Jahre & 10 Jahre & 15 Jahre\\
\midrule
\tabellenkoerper{data/flow_schwelle}
\bottomrule
\end{tabular}
\end{table}

\FloatBarrier

\vd{Die Tabelle beantwortet die Frage sch\"arfer als jedes Urteil. Im
mittleren Szenario liegt die Schwelle auf f\"unfzehn Jahre bei
\USDje{\FlowSchwelleBasis}; der heutige Kurs liegt
\Pct{\FlowAbstandSchwelle} dar\"uber. Das ist kein Bewertungsfehler, das ist
Abwesenheit von Sicherheitsabstand.

Entscheidend ist die obere Zeile. Schreibt man das zuletzt \emph{gemessene}
Wachstum fort, das des ersten Halbjahres 2026, liegt die Schwelle bei
\USDje{\FlowSchwelleUnten}, und der heutige Kurs liegt
\Pct{\FlowAbstandUnten} dar\"uber. Das Szenario, das die aktuellen Zahlen
tats\"achlich st\"utzen, ist also dasjenige, unter dem die Aktie deutlich zu
teuer ist. Genau darin liegt das Problem des Falls: Der Markt preist eine
Erholung ein, die in den Quartalsberichten noch nicht steht.}

\subsection{Ausl\"oser, die das Votum drehen}
\label{sec:ausloeser}

\vd{Beide sind kostenlos zu beobachten; einer wird quartalsweise
ver\"offentlicht, der andere steht jeden Tag im Kurszettel.

\emph{Zum Kauf.} Die organische Wachstumsrate wird \"uber zwei
aufeinanderfolgende Quartale wieder positiv. Dann tr\"agt das mittlere
Szenario, und der Kurs ist bis \USDje{\FlowSchwelleBFuenf} auf f\"unf Jahre
vertretbar. \emph{Oder} der Kurs f\"allt unter \USDje{\FlowSchwelleUnten},
ohne dass sich fundamental etwas verschlechtert; dann bekommt man selbst das
untere Szenario zum fairen Preis.

\emph{Zum Meiden.} Die organische Rate bleibt zwei weitere Quartale negativ.
Dann ist der Einbruch nicht mehr als Integrationsverdauung lesbar, sondern
als Kundenabgang aus der \"Ubernahme selbst, und es w\"are der teuerste
denkbare Verlauf: Der \Faktor{\KaufpreisJeUmsatz}-fache Jahresumsatz wurde
f\"ur genau diese Kunden bezahlt. Ebenfalls zum Meiden: eine
Wertberichtigung auf den Firmenwert von \MioUSD{\Goodwill}, oder ein erneuter
gro{\ss}er Zukauf, bevor der Verschuldungsgrad das Vor\"ubernahmeniveau von
\Faktor{\VerschuldungsgradVJ} wieder erreicht hat.}

\subsection{Was dieses Votum nicht wissen kann}

\vd{Die entscheidende Gr\"o{\ss}e ist nicht berechenbar: Wie wahrscheinlich
die Erholung des organischen Wachstums ist. Das ist eine Einsch\"atzung
\"uber Preiszyklen im Versicherungsmarkt und \"uber den Integrationserfolg,
und keine der ausgewerteten Prim\"arquellen tr\"agt sie.
Tabelle~\ref{tab:schwelle} sagt, was jede Annahme wert ist; welche zutrifft,
sagt sie nicht.

Zwei weitere Vorbehalte, die gegen das eigene Votum sprechen. Die Rechnung
schlie{\ss}t \emph{weitere Zuk\"aufe aus}, obwohl sie das Gesch\"aftsmodell
sind; funktioniert die Wiederanlage des freien Cash-Flows wie bisher, ist
\Pct{\FlowIRRBFuenfzehn} eher eine Untergrenze. Und der Kurs ist bereits
\Pct{\KursAbstandTief} \"uber seinem
Tief vom \Datum{\KurstagTiefZFVI} gelaufen; ein Teil dessen, worauf dieses
Votum warten will, hat der Markt vielleicht schon vorweggenommen.}
